\chapter{企业号关注流程及常见问题}

国科大的部分通知、身份认证和校园服务会通过微信企业号或企业微信相关入口完成。新生入学前后建议尽早完成关注和认证，以免错过后续办事信息。

\section{关注流程}

\begin{enumerate}
  \item 使用微信扫描学校发布的企业号认证入口。
  \item 按页面提示输入统一身份认证账号和密码，完成第一步认证。
  \item 补全认证信息。这里通常需要填写当前微信绑定的手机号或邮箱；手机号可在微信“设置 - 账号与安全 - 手机号”中查看。
  \item 认证成功后，返回微信通讯录或企业号入口确认是否能看到中国科学院大学企业号及相关应用。
  \item 若页面提示关注成功但应用未立即出现，可以退出微信稍后重进，或按下文方法排查。
\end{enumerate}

\section{手机号和微信号}

认证时填写的手机号应当与当前使用微信绑定的手机号一致。如果手机号不一致，系统可能显示认证成功，但企业号应用仍不可见。

如果需要查看或修改微信绑定手机号，可进入微信“设置 - 账号与安全 - 手机号”。如需查看或设置微信号，可在微信个人信息页中操作。

\section{解绑与重新认证}

如果已经关注过企业号，但需要重新认证，可以先进入身份认证相关页面，打开个人中心并退出当前企业号。若账号已经绑定微信，需要先解绑微信，再重新扫描学校发布的认证入口。

更换手机号时，可以在认证流程中根据页面提示修改手机号；若已经完成关注，则通常需要先解绑，再重新扫码认证。

\section{认证失败排查}

如果认证后只能看到“身份认证小助手”或“企业小助手”，看不到其他应用，可以依次排查：

\begin{enumerate}
  \item 退出微信客户端，稍后重新进入查看。认证状态有时会存在短暂延迟。
  \item 返回身份认证入口，重新输入认证时填写的手机号。
  \item 如果提示“该手机号不存在于企业通讯录中”，请核对自己是否扫了正确入口，并确认手机号是否与当前微信绑定手机号一致。
  \item 若仍无法解决，可取消关注后重新扫码认证。
  \item 如多次失败，可联系负责移动校园或本科生事务的老师，请老师在后台核对学工号、手机号和邮箱信息。
\end{enumerate}

\section{多手机号和多微信账号}

如果同学有多个手机号或多个微信账号，要特别注意认证时使用的是哪一个微信、哪一个手机号。常见问题是：曾经用手机号 A 认证过微信 A，后来又想用手机号 B 绑定的微信 B 认证；页面可能提示成功，但微信 B 仍看不到应用。

遇到这类情况时，建议先解绑原微信，再使用当前微信绑定的手机号重新认证。若原微信已经无法登录，需要请后台老师清除旧绑定关系后再重新认证。

\section{账号锁定与其他问题}

统一身份认证密码多次输入错误后，账号可能会临时锁定。若提示错误次数达到上限，可以等待一段时间后再试，或按学校统一身份认证系统的说明重置密码。

使用外国手机号认证时，国际区号前通常需要补两个 0。若页面提示手机号已存在，可能是该手机号曾绑定在其他学工号或旧账号下，需要由后台老师核对后处理。
